\begin{abstract}


% Large language models (LLMs) are increasingly used for tabular question answering, yet most evaluations assume clean, relational schemas. Real-world data, however, is often spread across multiple tables with non-relational presentations, such as hierarchical headers, merged cells, and cross-table dependencies. Manually curating such data is prohibitively expensive, and existing relational datasets often break when subjected to complex layout perturbations. To address this, we introduce \name, an automatic, from-scratch benchmark generator for evaluating numerical QA over multiple non-relational tables. \name separates supervision from presentation: it synthesizes a latent relational source to guarantee executable SQL ground-truth answers, but evaluates models exclusively on perturbed non-relational views. This design enables independent, contamination-free control over reasoning topology (parallel vs sequential), table count, layout noise, and unit heterogeneity. Our evaluation of recent LLMs across financial and environmental domains reveals a critical blind spot: non-relational formatting, unit discrepancies, and multi-table scaling cause sharp, compounded performance drops.

Large language models (LLMs) are increasingly used for tabular question answering, yet most evaluations assume clean, relational schemas. This leaves a gap for real-world analytical queries, where evidence may be distributed across multiple non-relational tables and encoded through layout, noise, units, or implicit cross-table links. Existing relational datasets are difficult to adapt to this setting because complex perturbations can break table validity and answer verification. 
We introduce \name, a from-scratch benchmark generator for controlled numerical QA evaluation over multiple non-relational tables. \name synthesizes a latent relational source for executable SQL supervision, then evaluates models on perturbed non-relational views derived from the same source. This separation preserves answer correctness while controlling reasoning topology, table count, layout perturbations, and unit heterogeneity. Our evaluation of LLMs across financial and environmental domains reveals a critical blind spot: non-relational formatting, unit discrepancies, and multi-table scaling cause sharp, compounded performance drops.
\end{abstract}